\IfFileExists{stacks-project.cls}{%
\documentclass{stacks-project}
}{%
\documentclass{amsart}
}

% For dealing with references we use the comment environment
\usepackage{verbatim}
\newenvironment{reference}{\comment}{\endcomment}
%\newenvironment{reference}{}{}
\newenvironment{slogan}{\comment}{\endcomment}
\newenvironment{history}{\comment}{\endcomment}

% For commutative diagrams we use Xy-pic
\usepackage[all]{xy}

% We use 2cell for 2-commutative diagrams.
\xyoption{2cell}
\UseAllTwocells

% We use multicol for the list of chapters between chapters
\usepackage{multicol}

% This is generally recommended for better output
\usepackage{lmodern}
\usepackage[T1]{fontenc}

% For cross-file-references
\usepackage{xr-hyper}

% Package for hypertext links:
\usepackage{hyperref}

% For any local file, say "hello.tex" you want to link to please
% use \externaldocument[hello-]{hello}
\externaldocument[introduction-]{introduction}
\externaldocument[conventions-]{conventions}
\externaldocument[sets-]{sets}
\externaldocument[categories-]{categories}
\externaldocument[topology-]{topology}
\externaldocument[sheaves-]{sheaves}
\externaldocument[sites-]{sites}
\externaldocument[stacks-]{stacks}
\externaldocument[fields-]{fields}
\externaldocument[algebra-]{algebra}
\externaldocument[brauer-]{brauer}
\externaldocument[homology-]{homology}
\externaldocument[derived-]{derived}
\externaldocument[simplicial-]{simplicial}
\externaldocument[more-algebra-]{more-algebra}
\externaldocument[smoothing-]{smoothing}
\externaldocument[modules-]{modules}
\externaldocument[sites-modules-]{sites-modules}
\externaldocument[injectives-]{injectives}
\externaldocument[cohomology-]{cohomology}
\externaldocument[sites-cohomology-]{sites-cohomology}
\externaldocument[dga-]{dga}
\externaldocument[dpa-]{dpa}
\externaldocument[sdga-]{sdga}
\externaldocument[hypercovering-]{hypercovering}
\externaldocument[schemes-]{schemes}
\externaldocument[constructions-]{constructions}
\externaldocument[properties-]{properties}
\externaldocument[morphisms-]{morphisms}
\externaldocument[coherent-]{coherent}
\externaldocument[divisors-]{divisors}
\externaldocument[limits-]{limits}
\externaldocument[varieties-]{varieties}
\externaldocument[topologies-]{topologies}
\externaldocument[descent-]{descent}
\externaldocument[perfect-]{perfect}
\externaldocument[more-morphisms-]{more-morphisms}
\externaldocument[flat-]{flat}
\externaldocument[groupoids-]{groupoids}
\externaldocument[more-groupoids-]{more-groupoids}
\externaldocument[etale-]{etale}
\externaldocument[chow-]{chow}
\externaldocument[intersection-]{intersection}
\externaldocument[pic-]{pic}
\externaldocument[weil-]{weil}
\externaldocument[adequate-]{adequate}
\externaldocument[dualizing-]{dualizing}
\externaldocument[duality-]{duality}
\externaldocument[discriminant-]{discriminant}
\externaldocument[derham-]{derham}
\externaldocument[local-cohomology-]{local-cohomology}
\externaldocument[algebraization-]{algebraization}
\externaldocument[curves-]{curves}
\externaldocument[resolve-]{resolve}
\externaldocument[models-]{models}
\externaldocument[functors-]{functors}
\externaldocument[equiv-]{equiv}
\externaldocument[pione-]{pione}
\externaldocument[etale-cohomology-]{etale-cohomology}
\externaldocument[proetale-]{proetale}
\externaldocument[relative-cycles-]{relative-cycles}
\externaldocument[more-etale-]{more-etale}
\externaldocument[trace-]{trace}
\externaldocument[crystalline-]{crystalline}
\externaldocument[spaces-]{spaces}
\externaldocument[spaces-properties-]{spaces-properties}
\externaldocument[spaces-morphisms-]{spaces-morphisms}
\externaldocument[decent-spaces-]{decent-spaces}
\externaldocument[spaces-cohomology-]{spaces-cohomology}
\externaldocument[spaces-limits-]{spaces-limits}
\externaldocument[spaces-divisors-]{spaces-divisors}
\externaldocument[spaces-over-fields-]{spaces-over-fields}
\externaldocument[spaces-topologies-]{spaces-topologies}
\externaldocument[spaces-descent-]{spaces-descent}
\externaldocument[spaces-perfect-]{spaces-perfect}
\externaldocument[spaces-more-morphisms-]{spaces-more-morphisms}
\externaldocument[spaces-flat-]{spaces-flat}
\externaldocument[spaces-groupoids-]{spaces-groupoids}
\externaldocument[spaces-more-groupoids-]{spaces-more-groupoids}
\externaldocument[bootstrap-]{bootstrap}
\externaldocument[spaces-pushouts-]{spaces-pushouts}
\externaldocument[spaces-chow-]{spaces-chow}
\externaldocument[groupoids-quotients-]{groupoids-quotients}
\externaldocument[spaces-more-cohomology-]{spaces-more-cohomology}
\externaldocument[spaces-simplicial-]{spaces-simplicial}
\externaldocument[spaces-duality-]{spaces-duality}
\externaldocument[formal-spaces-]{formal-spaces}
\externaldocument[restricted-]{restricted}
\externaldocument[spaces-resolve-]{spaces-resolve}
\externaldocument[formal-defos-]{formal-defos}
\externaldocument[defos-]{defos}
\externaldocument[cotangent-]{cotangent}
\externaldocument[examples-defos-]{examples-defos}
\externaldocument[algebraic-]{algebraic}
\externaldocument[examples-stacks-]{examples-stacks}
\externaldocument[stacks-sheaves-]{stacks-sheaves}
\externaldocument[criteria-]{criteria}
\externaldocument[artin-]{artin}
\externaldocument[quot-]{quot}
\externaldocument[stacks-properties-]{stacks-properties}
\externaldocument[stacks-morphisms-]{stacks-morphisms}
\externaldocument[stacks-limits-]{stacks-limits}
\externaldocument[stacks-cohomology-]{stacks-cohomology}
\externaldocument[stacks-perfect-]{stacks-perfect}
\externaldocument[stacks-introduction-]{stacks-introduction}
\externaldocument[stacks-more-morphisms-]{stacks-more-morphisms}
\externaldocument[stacks-geometry-]{stacks-geometry}
\externaldocument[moduli-]{moduli}
\externaldocument[moduli-curves-]{moduli-curves}
\externaldocument[examples-]{examples}
\externaldocument[exercises-]{exercises}
\externaldocument[guide-]{guide}
\externaldocument[desirables-]{desirables}
\externaldocument[coding-]{coding}
\externaldocument[obsolete-]{obsolete}
\externaldocument[fdl-]{fdl}
\externaldocument[index-]{index}

% Theorem environments.
%
\theoremstyle{plain}
\newtheorem{theorem}[subsection]{Theorem}
\newtheorem{proposition}[subsection]{Proposition}
\newtheorem{lemma}[subsection]{Lemma}

\theoremstyle{definition}
\newtheorem{definition}[subsection]{Definition}
\newtheorem{example}[subsection]{Example}
\newtheorem{exercise}[subsection]{Exercise}
\newtheorem{situation}[subsection]{Situation}

\theoremstyle{remark}
\newtheorem{remark}[subsection]{Remark}
\newtheorem{remarks}[subsection]{Remarks}

\numberwithin{equation}{subsection}

% Macros
%
\def\lim{\mathop{\mathrm{lim}}\nolimits}
\def\colim{\mathop{\mathrm{colim}}\nolimits}
\def\Spec{\mathop{\mathrm{Spec}}}
\def\Hom{\mathop{\mathrm{Hom}}\nolimits}
\def\Ext{\mathop{\mathrm{Ext}}\nolimits}
\def\SheafHom{\mathop{\mathcal{H}\!\mathit{om}}\nolimits}
\def\SheafExt{\mathop{\mathcal{E}\!\mathit{xt}}\nolimits}
\def\Sch{\mathit{Sch}}
\def\Mor{\mathop{\mathrm{Mor}}\nolimits}
\def\Ob{\mathop{\mathrm{Ob}}\nolimits}
\def\Sh{\mathop{\mathit{Sh}}\nolimits}
\def\NL{\mathop{N\!L}\nolimits}
\def\CH{\mathop{\mathrm{CH}}\nolimits}
\def\proetale{{pro\text{-}\acute{e}tale}}
\def\etale{{\acute{e}tale}}
\def\QCoh{\mathit{QCoh}}
\def\Ker{\mathop{\mathrm{Ker}}}
\def\Im{\mathop{\mathrm{Im}}}
\def\Coker{\mathop{\mathrm{Coker}}}
\def\Coim{\mathop{\mathrm{Coim}}}

% Boxtimes
%
\DeclareMathSymbol{\boxtimes}{\mathbin}{AMSa}{"02}

%
% Macros for moduli stacks/spaces
%
\def\QCohstack{\mathcal{QC}\!\mathit{oh}}
\def\Cohstack{\mathcal{C}\!\mathit{oh}}
\def\Spacesstack{\mathcal{S}\!\mathit{paces}}
\def\Quotfunctor{\mathrm{Quot}}
\def\Hilbfunctor{\mathrm{Hilb}}
\def\Curvesstack{\mathcal{C}\!\mathit{urves}}
\def\Polarizedstack{\mathcal{P}\!\mathit{olarized}}
\def\Complexesstack{\mathcal{C}\!\mathit{omplexes}}
% \Pic is the operator that assigns to X its picard group, usage \Pic(X)
% \Picardstack_{X/B} denotes the Picard stack of X over B
% \Picardfunctor_{X/B} denotes the Picard functor of X over B
\def\Pic{\mathop{\mathrm{Pic}}\nolimits}
\def\Picardstack{\mathcal{P}\!\mathit{ic}}
\def\Picardfunctor{\mathrm{Pic}}
\def\Deformationcategory{\mathcal{D}\!\mathit{ef}}
\usepackage{cleveref}
\crefname{originaldefinition}{original definition}{original definitions}
\Crefname{originaldefinition}{Original definition}{Original definitions}
\makeatletter
\newlabel{definition-cofinal}{{04E6}{source}{}{sourcepointer-1}{}}
\newlabel{definition-cofinal@cref}{{[originaldefinition][0][]04E6}{[1][0][]source}{}{}{}}
\makeatother
\crefname{originaldefinition}{original definition}{original definitions}
\Crefname{originaldefinition}{Original definition}{Original definitions}
\makeatletter
\newlabel{definition-system-over-poset}{{0030}{source}{}{sourcepointer-2}{}}
\newlabel{definition-system-over-poset@cref}{{[originaldefinition][0][]0030}{[1][0][]source}{}{}{}}
\makeatother
\crefname{originallemma}{original lemma}{original lemmas}
\Crefname{originallemma}{Original lemma}{Original lemmas}
\makeatletter
\newlabel{lemma-cofinal}{{04E7}{source}{}{sourcepointer-3}{}}
\newlabel{lemma-cofinal@cref}{{[originallemma][0][]04E7}{[1][0][]source}{}{}{}}
\makeatother
\crefname{originallemma}{original lemma}{original lemmas}
\Crefname{originallemma}{Original lemma}{Original lemmas}
\makeatletter
\newlabel{lemma-functorial-colimit}{{002K}{source}{}{sourcepointer-4}{}}
\newlabel{lemma-functorial-colimit@cref}{{[originallemma][0][]002K}{[1][0][]source}{}{}{}}
\makeatother
\crefname{originallemma}{original lemma}{original lemmas}
\Crefname{originallemma}{Original lemma}{Original lemmas}
\makeatletter
\newlabel{lemma-functorial-limit}{{002L}{source}{}{sourcepointer-5}{}}
\newlabel{lemma-functorial-limit@cref}{{[originallemma][0][]002L}{[1][0][]source}{}{}{}}
\makeatother
\title{Filtered index categories reduce to directed systems via a cofinal directed set}
\author{The Stacks Project Authors}
\date{stacks-project:categories.tex:local-snapshot-20261012; GNU FDL 1.2 or later}
\begin{document}
\maketitle
\noindent\textit{Editorial note.} Counted span is the author's own proof environment, source lines 2961-3063, following the statement at 2926-2959.\par\medskip
\noindent\textit{Editorial note.} Conclusion (2) is reduced by the author to (1) via the C\textasciicircum{}opp dictionary with details omitted (2973-2975); the reduction is definitional (I stays filtered) and is recorded here.\par\medskip
\noindent\textit{Editorial note.} Footnote at 2967-2969: the construction in fact produces a directed partially ordered set.\par\medskip
\noindent\textit{Allowed standard prerequisites (editorial).}\begin{itemize}
\item Colimits invariant under cofinal functors. Lemma lemma-cofinal of this chapter, invoked at proof lines 2969-2972 and 3019-3021. Transfers colimits from I to the constructed cofinal I0 and from I to I x omega via the projection Pi.
\item Functorial colimit comparison transformation. Lemma lemma-functorial-colimit of this chapter, referenced in the statement at 2938-2945. The transformation theta whose invertibility is part of conclusion (1); the cofinal functor identifies both colimits.
\item Cofinality of a functor. Definition definition-cofinal of this chapter, used at 3043-3061. Conditions (1) and (2) are verified for M0 at the end of the proof.
\end{itemize}\medskip
\ifdefined\hypertarget\hypertarget{sourcepointer-1}{}\fi
\noindent\textit{Source reference.} Original definition 04E6: sources/stacks-categories/pinned-original-source.tex, source lines 1943--1943. Navigation pointer only: the counted proof cites this statement by \textbackslash{}ref as a standing input; the deduction the counted proof has to make is carried out inside the counted span, and no part of that deduction is delegated to this pointer.\par\medskip
\ifdefined\hypertarget\hypertarget{sourcepointer-2}{}\fi
\noindent\textit{Source reference.} Original definition 0030: sources/stacks-categories/pinned-original-source.tex, source lines 2824--2824. Navigation pointer only: the counted proof cites this statement by \textbackslash{}ref as a standing input; the deduction the counted proof has to make is carried out inside the counted span, and no part of that deduction is delegated to this pointer.\par\medskip
\ifdefined\hypertarget\hypertarget{sourcepointer-3}{}\fi
\noindent\textit{Source reference.} Original lemma 04E7: sources/stacks-categories/pinned-original-source.tex, source lines 1977--1977. Navigation pointer only: the counted proof cites this statement by \textbackslash{}ref as a standing input; the deduction the counted proof has to make is carried out inside the counted span, and no part of that deduction is delegated to this pointer.\par\medskip
\ifdefined\hypertarget\hypertarget{sourcepointer-4}{}\fi
\noindent\textit{Source reference.} Original lemma 002K: sources/stacks-categories/pinned-original-source.tex, source lines 1654--1654. Navigation pointer only: the counted proof cites this statement by \textbackslash{}ref as a standing input; the deduction the counted proof has to make is carried out inside the counted span, and no part of that deduction is delegated to this pointer.\par\medskip
\ifdefined\hypertarget\hypertarget{sourcepointer-5}{}\fi
\noindent\textit{Source reference.} Original lemma 002L: sources/stacks-categories/pinned-original-source.tex, source lines 1688--1688. Navigation pointer only: the counted proof cites this statement by \textbackslash{}ref as a standing input; the deduction the counted proof has to make is carried out inside the counted span, and no part of that deduction is delegated to this pointer.\par\medskip
\begin{lemma}
\label{lemma-directed-category-system}
Let $\mathcal{I}$ be a filtered index category.
There exist a directed set $I$
and a system $(x_i, \varphi_{ii'})$ over $I$ in $\mathcal{I}$
with the following properties:
\begin{enumerate}
\item For every category $\mathcal{C}$ and every diagram
$M : \mathcal{I} \to \mathcal{C}$ with values in $\mathcal{C}$,
denote $(M(x_i), M(\varphi_{ii'}))$
the corresponding system over $I$. If
$\colim_{i \in I} M(x_i)$ exists then so does
$\colim_\mathcal{I} M$ and the transformation
$$
\theta :
\colim_{i \in I} M(x_i)
\longrightarrow
\colim_\mathcal{I} M
$$
of Lemma \ref{lemma-functorial-colimit} is an isomorphism.
\item For every category $\mathcal{C}$ and every diagram
$M : \mathcal{I}^{opp} \to \mathcal{C}$ in $\mathcal{C}$, denote
$(M(x_i), M(\varphi_{ii'}))$ the corresponding inverse system
over $I$. If $\lim_{i \in I} M(x_i)$ exists then so does
$\lim_{\mathcal{I}^{opp}} M$ and the transformation
$$
\theta :
\lim_{\mathcal{I}^{opp}} M
\longrightarrow
\lim_{i \in I} M(x_i)
$$
of Lemma \ref{lemma-functorial-limit} is an isomorphism.
\end{enumerate}
\end{lemma}

\begin{proof}
As explained in the text following
Definition \ref{definition-system-over-poset}, we may view
preordered sets as categories and systems as functors.
Throughout the proof, we will freely shift between these two points of view.
We prove the first statement by constructing a category
$\mathcal{I}_0$, corresponding to a directed set\footnote{In fact,
our construction will produce a directed partially ordered set.}, and a cofinal
functor $M_0 : \mathcal{I}_0 \to \mathcal{I}$. Then, by
Lemma \ref{lemma-cofinal}, the colimit of a diagram
$M : \mathcal{I} \to \mathcal{C}$ coincides with the
colimit of the diagram $M \circ M_0 : \mathcal{I}_0 \to \mathcal{C}$,
from which the statement follows. The second statement is dual to the
first and may be proved by interpreting a limit in $\mathcal{C}$ as
a colimit in $\mathcal{C}^{opp}$. We omit the details.

\medskip\noindent
A category $\mathcal{F}$ is called {\em finitely generated} if
there exists a finite set $F$ of arrows in $\mathcal{F}$, such that
each arrow in $\mathcal{F}$ may be obtained by composing
arrows from $F$. In particular, this implies that $\mathcal{F}$ has
finitely many objects. We start the proof by reducing to the case
when $\mathcal{I}$ has the property that every finitely generated
subcategory of $\mathcal{I}$ may be extended to a finitely
generated subcategory with a unique final object.

\medskip\noindent
Let $\omega$ denote the directed set of finite ordinals, which
we view as a filtered category. It is easy to verify that the
product category $\mathcal{I}\times \omega$ is also filtered,
and the projection
$\Pi : \mathcal{I} \times \omega \to \mathcal{I}$
is cofinal.

\medskip\noindent
Now let $\mathcal{F}$ be any finitely generated
subcategory of $\mathcal{I}\times \omega$.
By using the axioms of a filtered category and a simple induction
argument on a finite set of generators of $\mathcal{F}$,
we may construct a cocone $(\{f_i\}, i_\infty)$ in $\mathcal{I} \times \omega$
for the diagram $\mathcal{F} \to \mathcal{I} \times \omega$. That is, a morphism
$f_i : i \to i_\infty$ for every object $i$ in $\mathcal{F}$
such that for each arrow $f : i \to i'$ in $\mathcal{F}$
we have $f_i = f_{i'} \circ f$. We can also choose $i_\infty$ such
that there are no arrows from $i_\infty$ to an object in $\mathcal{F}$.
This is possible since
we may always post-compose the arrows $f_i$ with an arrow
which is the identity on the $\mathcal{I}$-component and
strictly increasing on the $\omega$-component.
Now let $\mathcal{F}^+$ denote the category consisting of all
objects and arrows in $\mathcal{F}$
together with the object $i_\infty$, the identity
arrow $\text{id}_{i_\infty}$ and the arrows $f_i$.
Since there are no arrows from $i_\infty$ in $\mathcal{F}^+$
to any object of $\mathcal{F}$, the arrow set in $\mathcal{F}^+$
is closed under composition, so $\mathcal{F}^+$ is indeed
a category. By construction, it is a finitely
generated subcategory of $\mathcal{I}$ which has $i_\infty$ as
unique final object. Since, by Lemma \ref{lemma-cofinal},
the colimit of any diagram $M : \mathcal{I} \to \mathcal{C}$
coincides with the colimit of $M\circ\Pi$ , this gives the desired
reduction.

\medskip\noindent
The set of all finitely generated subcategories of $\mathcal{I}$
with a unique final object is naturally ordered by inclusion.
We take $\mathcal{I}_0$ to be the category corresponding
to this set. We also have a functor
$M_0 : \mathcal{I}_0 \to \mathcal{I}$, which takes an
arrow $\mathcal{F} \subset \mathcal{F'}$ in
$\mathcal{I}_0$ to the unique map from the final object of
$\mathcal{F}$ to the final object of $\mathcal{F}'$.
Given any two finitely generated subcategories of
$\mathcal{I}$, the category generated by these two categories is
also finitely generated. By our assumption on $\mathcal{I}$, it is
also contained in a finitely generated subcategory of $\mathcal{I}$
with a unique final object. This shows that $\mathcal{I}_0$ is directed.

\medskip\noindent
Finally, we verify that $M_0$ is cofinal. Since any
object of $\mathcal{I}$ is the final object in the subcategory
consisting of only that object and its identity arrow, the functor
$M_0$ is surjective on objects. In particular, Condition (1) of
Definition \ref{definition-cofinal} is satisfied. Given
an object $i$ of $\mathcal{I}$, objects $\mathcal{F}_1, \mathcal{F}_2$ in
$\mathcal{I}_0$ and maps $\varphi_1 : i \to M_0(\mathcal{F}_1)$
and $\varphi_2 : i \to M_0(\mathcal{F}_2)$ in
$\mathcal{I}$, we can take $\mathcal{F}_{12}$ to be a finitely
generated category with a unique final object containing
$\mathcal{F}_1$, $\mathcal{F}_2$ and the morphisms $\varphi_1, \varphi_2$.
The resulting diagram commutes
$$
\xymatrix{
& M_0(\mathcal{F}_{12}) & \\
M_0(\mathcal{F}_{1}) \ar[ru] & & M_0(\mathcal{F}_{2}) \ar[lu] \\
& i \ar[lu] \ar[ru]
}
$$
since it lives in the category $\mathcal{F}_{12}$ and
$M_0(\mathcal{F}_{12})$ is final in
this category. Hence also Condition (2) is satisfied, which concludes
the proof.
\end{proof}

\end{document}
